\section{평가}

\subsection{평가 설계 및 질문}

본 평가는 세 가지 독립적인 측면에서 제안 시스템을 검증한다.
평가~A는 PQC 도구 체인 통합 과정에서 발생하는 인프라 수준 함정을 카탈로그화하고, 표준 보안 도구의 사전 탐지 가능 여부를 분석한다.
평가~B는 LLM 기반 레거시 코드 마이그레이션의 정확도를 정량적으로 측정하고, 검증 계층별 역할을 분리 분석한다.
평가~C는 제안 파이프라인이 평가~A·B에서 확인된 함정에 어떤 차단·감지 효과를 가지는지 측정한다.

각 평가에서 답하는 질문은 다음과 같다.

\begin{itemize}
  \item \textbf{Q1} (평가~A): PQC 도구 체인 통합 시 어떤 종류의 함정이 발생하며, 표준 SCA/SAST/SBOM 도구는 이를 정적·SBOM 기반으로 사전 탐지할 수 있는가?
  \item \textbf{Q2} (평가~B): LLM 기반 RSA\textrightarrow ML-KEM/ML-DSA 코드 마이그레이션은 어느 검증 계층에서 실패하며, 정적 AST 검증과 의미적 실행 검증은 어떤 실패 유형을 각각 차단하는가?
  \item \textbf{Q3} (평가~C): 제안 파이프라인은 표준 도구가 사전 탐지하지 못하는 함정 6종을 얼마나 차단·감지하는가?
  \item \textbf{Q4} (평가~C): LLM 단독 대비 다층 검증 적용 시 실패 코드의 production 전달 위험은 얼마나 감소하는가?
\end{itemize}

\subsection{실험 환경}

\textbf{파이프라인 실행 환경}은 GitHub Actions ubuntu-latest 러너를 사용한다.
OQS-Nginx 서버는 OpenSSL~3.4.0, liboqs~0.13.0, oqs-provider~0.9.0, nginx~1.27.2 소스 빌드 기반이며, TLS 검증 클라이언트는 OQS-curl(\texttt{openquantumsafe/curl:0.11.0})을 사용한다.

\textbf{LLM 마이그레이션 평가 환경}은 Windows~10 Pro, Docker Desktop, \texttt{python:3.11-slim} 기반 컨테이너에서 수행하였다.
liboqs~0.14.0, liboqs-python~0.14.1을 사용하며, LLM 엔드포인트는 GitHub Models API(\texttt{gpt-4o-mini}, temperature=0)다.
측정 시점은 2026-05-08이며, 총 80 trial(16 패턴 $\times$ k=5)을 실행하였다.

\begin{table}[htbp]
\centering
\renewcommand{\arraystretch}{1.3}
\caption{실험 환경 및 도구 버전}
\label{tab:env}
\begin{tabular}{ll}
\toprule
구성 요소 & 버전 \\
\midrule
OpenSSL & 3.4.0 (소스 빌드) \\
liboqs (서버) & 0.13.0 \\
liboqs (평가 하네스) & 0.14.0 \\
oqs-provider & 0.9.0 \\
nginx & 1.27.2 \\
Trivy & 0.69.3 \\
Semgrep & 1.156.0 \\
Gitleaks & 8.21.2 \\
LLM 모델 & gpt-4o-mini (temperature=0) \\
\bottomrule
\end{tabular}
\end{table}

\subsection{평가 A: PQC 도구 체인 함정 카탈로그}

본 연구의 OQS 스택 통합 과정에서 발생한 인프라 수준 함정 6종을 트리거·증상·원인·방어책 단위로 구조화하였다.

\begin{table}[htbp]
\centering
\renewcommand{\arraystretch}{1.3}
\caption{PQC 도구 체인 함정 6종 카탈로그}
\label{tab:trap_catalog}
\begin{tabular}{p{0.5cm}p{3cm}p{4.5cm}p{4cm}}
\toprule
ID & 함정 & 증상 & 방어책 \\
\midrule
F-1 & oqs-provider 버전 비대칭 & TLS 1.3 핸드셰이크 \texttt{alert 552} & 자체 빌드로 라이브러리 라인 통일 \\
F-2 & OpenSSL 1.1.1 $\leftrightarrow$ 3.x OID 불일치 & \texttt{SSL\_CTX\_use\_certificate failed} & cert-builder를 nginx와 동일 라인으로 통일 \\
F-3 & ML-KEM IETF code point 충돌 & nginx 파싱 통과, 핸드셰이크 group 매칭 0 & 동일 라이브러리 라인 통일 \\
F-4 & PQ key strength 0 보고 & \texttt{ee key too small (SECLEVEL)} & \texttt{@SECLEVEL=0} + TLSv1.3 + group whitelist \\
F-5 & Docker 이미지 태그/빌드 재현성 문제 & 이전 CI 통과 이미지 조합이 이후 재구성 시 실패 & 자체 빌드로 외부 이미지 태그 의존 제거 \\
F-6 & standalone PQ group 미협상 & nginx 기동 정상, 실제 핸드셰이크 실패 & hybrid group으로 회귀 + 매트릭스 CI 감지 \\
\bottomrule
\end{tabular}
\end{table}

표준 보안 도구의 사전 탐지 가능 여부를 분석한 결과, Trivy(SCA), Semgrep(SAST), CycloneDX CBOM 중 어느 도구도 6종 함정을 정적·SBOM 기반으로 사전 탐지하지 못하였다(0/6).
이는 표준 도구가 알려진 CVE·설정 오류·암호 자산 목록화에 집중하며, OQS 스택의 라이브러리 라인 비대칭이나 code point 불일치와 같은 통합 수준 함정은 사전 탐지 범위 밖임을 의미한다.

\subsection{평가 B: LLM 기반 코드 마이그레이션 정확도}

\subsubsection{평가 설계}

입력 패턴은 라이브러리(cryptography.hazmat, PyCryptodome) $\times$ 사용 케이스(keygen, encrypt, sign, kex) $\times$ 복잡도(simple, complex)의 직교 설계 N=16으로 구성하였다.
각 패턴에 대해 RSA에서 ML-KEM-768/ML-DSA-65(oqs-python) 기반 코드로의 변환을 gpt-4o-mini에 요청하고, 4단계 검증 계층으로 결과를 측정하였다.

4단계 검증 계층(L1\textasciitilde L4)은 다음과 같이 정의한다.
L1은 LLM API 응답 수신 여부, L2는 AST 기반 정적 검증(oqs API 사용 여부, 인스턴스 생성 패턴, 클래스 직접 호출 금지), L3는 패턴별 round-trip 의미적 실행 검증, L4는 reference와 candidate의 AST feature 가중 평균 동등성 검증(임계값 0.8)이다.

하네스 sanity check로 정답 16개를 candidate로 입력 시 L2/L3/L4 모두 16/16 통과, 원본 RSA 코드 입력 시 0/16으로 false positive/negative 0건을 확인하였다.

\subsubsection{다층 ablation 결과}

\begin{table}[htbp]
\centering
\renewcommand{\arraystretch}{1.3}
\caption{4단계 검증 계층별 통과율 (총 80 trial)}
\label{tab:llm_ablation}
\begin{tabular}{llll}
\toprule
계층 & 통과 / 80 & 전체 비율 & 응답 55건 기준 \\
\midrule
L1 (LLM 응답 수신) & 55 & 68.8\% & 100\% \\
L2 (AST 검증) & 24 & 30.0\% & 43.6\% \\
L3 (의미적 실행 검증) & 0 & 0.0\% & 0.0\% \\
L4 (동등성 검증) & 40 & 50.0\% & 72.7\% \\
\bottomrule
\end{tabular}
\end{table}

응답을 받은 55 trial 모두 어딘가에서 실패하여 의미적으로 정확한 코드를 생산한 경우는 0건이었다.
L4 통과율이 L3보다 높다는 점은 AST feature 수준의 구조적 유사성이 실제 실행 가능성을 보장하지 않음을 보여준다.

\subsubsection{실패 모드 taxonomy}

\begin{table}[htbp]
\centering
\renewcommand{\arraystretch}{1.3}
\caption{LLM 실패 모드 7종 및 차단 계층}
\label{tab:llm_taxonomy}
\begin{tabular}{p{5cm}ll}
\toprule
실패 모드 & 빈도 & 차단 계층 \\
\midrule
missing\_instance\_construction & 20 (36.4\%) & L2 \\
wrong\_arg\_count & 16 (29.1\%) & L3 \\
no\_migration & 6 (10.9\%) & L2 \\
class\_direct\_call & 5 (9.1\%) & L2 \\
wrong\_kwarg & 5 (9.1\%) & L3 \\
method\_hallucination & 2 (3.6\%) & L3 \\
algorithm\_name\_hallucination & 1 (1.8\%) & L3 \\
\bottomrule
\end{tabular}
\end{table}

L2와 L3가 차단하는 실패 모드 카테고리의 overlap은 0이다.
\texttt{wrong\_arg\_count}, \texttt{wrong\_kwarg}, \texttt{method\_hallucination}, \texttt{algorithm\_name\_hallucination}은 100\% L3에서만, \texttt{missing\_instance\_construction}, \texttt{no\_migration}, \texttt{class\_direct\_call}은 100\% L2에서 차단된다.
따라서 정적 AST 검증 단독으로는 시그니처 환각 24건(44\%)을 차단하지 못하며, 의미적 실행 검증을 추가해야 이를 차단할 수 있다.

\subsection{평가 C: 파이프라인의 함정 차단 및 감지 효과}

\subsubsection{인프라 함정에 대한 파이프라인 효과}

\begin{table}[htbp]
\centering
\renewcommand{\arraystretch}{1.3}
\caption{함정 6종에 대한 파이프라인 layer별 효과}
\label{tab:pipeline_effect}
\begin{tabular}{p{0.5cm}p{3cm}ccc}
\toprule
ID & 함정 & 표준 도구 & 자체 빌드 & 매트릭스 CI \\
\midrule
F-1 & sigalg drift & $\times$ & $\checkmark$ 예방 & $\checkmark$ 감지 \\
F-2 & OID mismatch & $\times$ & $\checkmark$ 예방 & $\checkmark$ 감지 \\
F-3 & group code point & $\times$ & $\checkmark$ 예방 & $\checkmark$ 감지 \\
F-4 & SECLEVEL 거부 & $\times$ & $\checkmark$ 예방 & $\checkmark$ 감지 \\
F-5 & 이미지 태그 재현성 & $\times$ & $\checkmark$ 예방 & $\checkmark$ 감지 \\
F-6 & standalone 미협상 & $\times$ & --- & $\checkmark$ 감지 \\
\midrule
\textbf{합계} & & \textbf{0/6} & \textbf{5/6} & \textbf{6/6} \\
\bottomrule
\end{tabular}
\end{table}

자체 빌드는 6종 함정 중 5개를 사전 예방하고, 매트릭스 CI는 6종 모두를 런타임에서 감지한다.
표준 SCA/SAST/SBOM은 6종 어느 것도 정적·SBOM 기반으로 사전 탐지하지 못한다(0/6).
F-6(standalone PQ group 미협상)은 자체 빌드로도 예방 불가이며 hybrid group 회귀와 매트릭스 CI 감지로 대응한다.

\subsubsection{LLM 다층 방어 효과}

\begin{table}[htbp]
\centering
\renewcommand{\arraystretch}{1.3}
\caption{검증 계층 적용 시나리오별 실패 코드 전달 위험}
\label{tab:crash_risk}
\begin{tabular}{lll}
\toprule
시나리오 & 통과 trial & 통과 코드의 런타임 실패 위험 \\
\midrule
① 검증 없음 (LLM 출력 직접 적용) & 55 & 100\% \\
② AST 검증만 (L2) & 24 & 100\% \\
③ AST + 의미적 검증 (L2 + L3) & 0 & 해당 없음 \\
\bottomrule
\end{tabular}
\end{table}

AST와 의미적 실행 검증을 모두 적용한 경우, 본 평가에서는 production 단계로 전달되는 코드가 0건이 되어 실패 코드의 전달을 차단하였다.
이는 LLM 마이그레이션이 성공적으로 완료되었음을 의미하는 것이 아니라, 현재 모델 출력이 의미적 검증을 통과하지 못해 위험 코드가 배포 단계로 전달되지 않았음을 의미한다.

\subsubsection{성능 검증 결과}

Stage~2 환경에서 TLS 핸드셰이크 시간을 30회 측정한 결과, 중앙값 44.575ms, 실패율 0\%로 절대 임계치(3,000ms)를 통과하였다.
상대 임계치 판정은 Stage~1 baseline 아티팩트를 기준으로 수행되도록 설계되며, 본 실험에서는 강제 실패 조건(\texttt{PERF\_FORCE\_FAIL=1})을 통해 \texttt{rollback.sh}가 롤백 전용 브랜치와 PR을 생성하는 경로를 검증하였다.

\subsection{논의}

평가~A와 평가~B는 PQC 전환 과정의 함정이 인프라 통합 수준과 코드 마이그레이션 수준에서 각각 발생하며, 두 층 모두 표준 도구 단독으로는 대응하기 어렵다는 것을 보여준다.
평가~C는 자체 빌드(5/6 예방)와 매트릭스 CI(6/6 감지)의 조합이 인프라 함정에 대한 표준 도구의 사각지대를 메우고, AST + 의미적 검증의 다층 구조가 실패 코드의 production 전달을 차단함을 정량적으로 입증한다.

본 평가의 범위와 해석상 유의점은 다음과 같다.
인프라 함정 카탈로그는 N=1 프로젝트의 경험적 사례이므로 통계적 일반화에는 한계가 있다.
LLM 평가는 단일 모델(gpt-4o-mini)과 합성 패턴 N=16만을 대상으로 하며, 더 큰 모델이나 실제 OSS 코드에서의 결과는 다를 수 있다.
성능 측정은 GitHub Actions 러너와 Docker 환경에서 수행하였으므로 절대 성능값보다 동일 환경 내 상대 비교로 해석해야 한다.
상세 한계 및 향후 연구 방향은 5장에서 논의한다.